\uuid{9S9O}
\exo7id{3361}
\titre{exo7 3361}
\auteur{quercia}
\organisation{exo7}
\datecreate{2010-03-09}
\isIndication{false}
\isCorrection{false}
\chapitre{Matrice}
\sousChapitre{Inverse, méthode de Gauss}
\cpge{mpsi-mp2i.determinants.d, pcsi.calcul-matriciel-systemes-lineaires.d, ptsi.calcul-matriciel-systemes-lineaires.d, mpsi-mp2i.matrices-applications-lineaires.b, pcsi.matrices-applications-lineaires.b, ptsi.matrices-applications-lineaires.b, bcpst1.matrices}
\module{Algèbre}
\niveau{L1}
\difficulte{}

\contenu{
\question{
Soit $M \in \mathcal{M}_n(\Q)$. Comparer les énoncés :
\par {\bf1:} $M$ est inversible dans $\mathcal{M}_n(\Q)$.
\par {\bf2:} $M$ est inversible dans $\mathcal{M}_n(\C)$.
}
}
